\section{Methodology}
\label{sec:methodology}

The methodology proceeds from diagnosis to supervised testing. It begins with
a completed pilot (Section~\ref{sec:pilot}) that identifies errors in
fluent-only parsing, proceeds to corpus construction (Section~\ref{sec:data})
that converts the formal account into paired supervision, and ends with a model
experiment (Section~\ref{sec:model}) comparing normalising and repair-aware
targets.

\subsection{Pilot Study (Completed)}
\label{sec:pilot}

The completed pilot tested whether a Gemma-2 9B semantic parser
\citep{gemmateamGemma2Improving2024}, fine-tuned with LoRA only on fluent
sentence--SBN pairs \citep{huLoRALowRankAdaptation2021}, could recover the
fluent meaning from an input containing RSR. Following
\citet{senSemanticParsingDisfluent2021}, the RSR variants changed only the
natural-language input and were evaluated against the original fluent gold
SBN. Appendix~\ref{app:pilot-training}
reports the pilot training configuration.

\subsubsection{Parsing Results and Output Diagnosis}

The materials came from the English train-gold portion of PMB 5.1.0
\citep{abzianidzeParallelMeaningBank2017}. Because the source sentences and
their fluent SBN targets were included in the LoRA fine-tuning data, the parser
had been trained on the underlying fluent sentence--SBN mappings; only the
RSR-modified inputs were new. Qwen3-32B
\citep{yangQwen3TechnicalReport2025} selected a repair word and inserted a
different word of the same part of speech immediately before it as the
reparandum. The repair remained in place, keeping the intended meaning aligned
with the fluent gold SBN. After three distinct variants were generated for each source
and filtered, 833 sentences yielded 2,499 unmarked RSR inputs. Adding
\textit{I mean} between each reparandum and repair produced another 2,499.
Appendix~\ref{app:pilot-results} gives the full condition-level
results.

Each output was converted into a graph and compared with the fluent reference
SBN. The \emph{scorable rate} records how many outputs could be converted and
scored. For those outputs, SMATCH++ measures concept and relation overlap with
the reference \citep{opitzSMATCHStandardizedExtended2023}.
\emph{Failure-penalised F1} assigns zero to unscoreable outputs before averaging,
so graph-generation failures remain visible.

\begin{table}[htbp]
\centering
\small
\begin{tabular}{lrr}
\toprule
Condition & Scorable (\%) & Failure-penalised F1 \\
\midrule
Fluent      & 97.5 & 0.950$^{***}$ \\
Unmarked RSR & 94.6 & 0.786$^{***}$ \\
RSR + editing expression & 76.5 & 0.609$^{***}$ \\
\bottomrule
\end{tabular}
\caption{Parser performance across the Fluent, unmarked-RSR and
RSR-with-editing-expression input conditions. Asterisks indicate that all
three pairwise Wilcoxon signed-rank comparisons of failure-penalised F1 gave
$p<.001$.}
\label{tab:pilot}
\end{table}

Table~\ref{tab:pilot} shows a stepwise decline in both measures.
RSR lowered F1 while leaving most outputs scorable; adding \textit{I mean}
produced a further loss driven mainly by graph-generation failures. Among
outputs that remained scorable, the additional decline in F1 was much smaller.

The output patterns point to a common source of this decline. The parser often
mapped the words in the input into SBN without distinguishing their roles in
the repair. Of the 2,363 scorable RSR outputs, 90.1\% contained more concepts
than the parser's output for the corresponding fluent source sentence, with a median
of two additional concepts. Much of the fluent proposition therefore remained
recoverable, but the reparandum was often added to the asserted content rather
than interpreted as withdrawn.

Including \textit{I mean} introduced further errors into this surface-based
mapping.
Figure~\ref{fig:pilot-output} shows a parse that recovers the accepted
proposition about a cow giving birth, but also asserts the \textit{goat} and
interprets \textit{I mean} as a meaning event involving the speaker. No
structure records that the cow replaces the goat. The SBN is formally
possible, but it does not express the intended RSR.



\begin{figure}[htbp]
\centering
\small\textbf{Input:} A goat, I mean, a cow gave birth to a calf with two heads.
\par\smallskip
\begin{minipage}[t]{0.485\textwidth}
\centering
\footnotesize\textbf{Fluent reference SBN}
\begin{lstlisting}[style=sbn,basicstyle=\ttfamily\footnotesize,columns=fullflexible,breaklines=true,breakatwhitespace=true,breakindent=1em,xleftmargin=0pt,xrightmargin=0pt,framexleftmargin=2pt,framexrightmargin=2pt,aboveskip=0.25em,belowskip=0.2em]
cow.n.01    % A cow
...
\end{lstlisting}
\end{minipage}\hfill
\begin{minipage}[t]{0.485\textwidth}
\centering
\footnotesize\textbf{RSR + editing-expression parse}
\begin{lstlisting}[style=sbn,basicstyle=\ttfamily\footnotesize,columns=fullflexible,breaklines=true,breakatwhitespace=true,breakindent=1em,xleftmargin=0pt,xrightmargin=0pt,framexleftmargin=2pt,framexrightmargin=2pt,aboveskip=0.25em,belowskip=0.2em,escapeinside={(*@}{@*)}]
(*@\colorbox{yellow!30}{\strut\texttt{goat.n.01}}@*)    % A goat
(*@\colorbox{yellow!30}{\strut\texttt{person.n.01 EQU speaker}}@*)      %I
(*@\colorbox{yellow!30}{\strut\texttt{mean.v.01 Theme -2 Agent -1 Time +1}}@*)     % mean
(*@\colorbox{yellow!30}{\strut\texttt{time.n.08 EQU now}}@*)
cow.n.01    % A cow
...
\end{lstlisting}
\end{minipage}
\caption{Excerpts from the fluent reference and RSR parse. Highlighting marks
content incorrectly asserted from the reparandum and editing expression; shared
content is omitted.}
\label{fig:pilot-output}
\end{figure}

The output still receives a SMATCH++ score of 0.764 because the accepted
proposition continues to receive credit while the added content only lowers
the overlap. A substantial graph score can therefore conceal the semantic
error central to RSR. Appendix~\ref{app:pilot-goat-cow} gives the complete
RSR and fluent-input outputs.

\subsubsection{Recognition Probe and Implications}

The parsing results do not show whether the model failed to recognise RSR or
lacked an SBN pattern for representing it. I therefore removed the LoRA
adapter and asked the instruction-tuned Gemma-2 9B to mark the reparandum,
editing expression and repair in plain text. The probe detected 100.0\% of inputs
with an editing expression and 44.8\% without one, while falsely identifying RSR
in 1.0\% of fluent inputs. The editing expression was therefore a strong cue,
although detection did not ensure accurate segmentation.

The two surface forms exposed different weaknesses. Without an editing expression,
the base model missed more than half of the RSRs, while the parser often folded
both candidates into a fluent-looking graph. With \textit{I mean}, the probe
detected every RSR, but the parser often represented the expression as semantic
content. Because the probe changed both the task and the adaptation state, it
cannot isolate SBN generation as the source of error. The pilot also covered
one model size and included no repair-aware supervision. The main experiment
will therefore test repair-aware training across model sizes; the next section
describes its targets.

%==============NEW DATA SECTION
\subsection{Repair-Aware Corpus Construction}
\label{sec:data}

\subsubsection{Paired Construction}

No existing corpus pairs replacement self-repair utterances with SBN targets
that distinguish withdrawn from accepted content. The PMB contains fluent
sentence--SBN pairs \citep{abzianidzeParallelMeaningBank2017}, while speech and
disfluency resources contain self-repairs without corresponding SBN annotations
\citep{deshmukhResegmentationSWITCHBOARD1998,
passaliLARDLargescaleArtificial2022,guptaDisflQABenchmarkDataset2021}. Each PMB
source pair is therefore transformed in parallel into an RSR utterance and the
repair-aware SBN defined in Section~\ref{sec:paired-configuration}. Unlike the
pilot, which altered only the utterance, this construction changes both sides
of the pair.

Figure~\ref{fig:paired-transformation} summarises this parallel construction.

\begin{figure}[ht]
\centering
\begin{tikzcd}[column sep=8em, row sep=normal]
\text{fluent sentence} \arrow[r, "\text{insert reparandum}"] \arrow[d, "\text{PMB gold}"'] &
\text{RSR utterance} \arrow[d, dashed, "?"] \\
\text{fluent gold SBN} \arrow[r, "\text{RSR transformation}"'] &
\text{repair-aware SBN}
\end{tikzcd}
\caption{Parallel construction of an RSR utterance and repair-aware target.}
\label{fig:paired-transformation}
\end{figure}

Construction will start from a fluent sentence, its fluent gold SBN and its word--concept
alignment. An aligned noun, verb or adjective is selected as the repair, and a
different expression of the same part of speech is drawn from WordNet
\citep{millerWordNetLexicalDatabase1994} and inserted immediately before it.
For example, \textit{``She prefers tea''} becomes \textit{``She prefers
coffee, tea''}. The implemented corpus transformation replaces one aligned
lexical concept at a time; the repair follows the inserted reparandum either
immediately or after an intervening editing expression.

The repair keeps its gold representation. The transformation adds the
reparandum in the paired \sbnterm{CORRECTION} and \sbnterm{CONJUNCTION}
contexts, using the direct construction where possible and the anchor or
movement construction where required. It preserves the source graph
outside the selected replacement site. Each record stores the original and
repair-aware SBNs, the two replacement concepts, their status, their links to
shared content and the construction used. These fields pair each RSR input
with either target in the model comparison and record the information needed
to compute the structural scores in Section~\ref{sec:model}.
Appendix~\ref{app:generated-record}
shows a complete record.

\subsubsection{Validation and Selection}

Structural validation checks that the target forms a well-defined graph,
contains the paired RSR contexts, assigns withdrawn and accepted status
correctly and preserves the source graph elsewhere. It also checks the links
that connect both replacement concepts to shared content. A target that fails
any check is excluded.

A legal graph can still contain a poor training example. The competing
expressions must differ in meaning, yet both must fit the sentence context.
Bidirectional NLI scores identify replacements that remain mutually entailing
\citep{laurerLessAnnotatingMore2024}, while a causal language model measures
contextual fit as the mean token surprisal of the reparandum conditioned on
its left context. The NLI boundaries and the upper surprisal threshold are set
on a manually labelled calibration sample and frozen before the corpus is filtered. A separate sample
from the retained data estimates the proportion that satisfies both criteria.
Appendix~\ref{app:corpus-selection} gives the calibration and audit procedure.

\subsubsection{Experimental Partitions}

Every generated item inherits the train, development or test split of its PMB
source document. All variants of a source remain together, preventing a
training document from reappearing in evaluation. Every model condition also
retains the same fluent PMB component. Two training conditions receive
identical RSR utterances. The \emph{normalising} condition pairs each utterance
with its fluent gold SBN, whereas the \emph{repair-aware} condition pairs it
with the generated SBN that retains the replacement structure.

% 2026-10-24: additive recipe supersedes the original one-variant rule.
Every base RSR training item is retained without an editing expression.
An additional marked variant is constructed for approximately 30\% of base
items using one of four editing expressions---\textit{I mean}, \textit{no},
\textit{no wait} or \textit{actually}---with the same SBN target.
The current set contains 3,489 base items and 1,047 additional marked
variants, for 4,536 RSR training rows; marked variants account for 23.1\%
of these rows. Thus training retains the fluent component, adds base RSR,
and then adds marked variants. Selected items have greater row weight;
the normalising and repair-aware conditions use identical inputs and sampling.
Each test item contributes
two matched utterances with the same repair content and target.
One contains no editing expression, while the other contains one. The test
expressions include both forms seen during training and the held-out
\textit{sorry}, allowing the effect of an editing expression and
generalisation to an unseen expression to be evaluated separately.

Two test-only challenge sets contain adversarial examples that examine
behaviour that positive, sentence-local RSR examples cannot establish. The
non-RSR set contains lists, appositions and similar fluent sequences that
exhibit local surface patterns resembling RSR without expressing replacement.
It tests whether the parser distinguishes genuine repair from superficial
similarity; assigning an RSR configuration to these items counts as a false
positive. The discourse set contains controls without RSR, later reference to
the accepted referent, reference to both replacement candidates and
corrections between discourse turns. The last condition tests whether the
within-clause configuration is applied outside its scope.
Appendix~\ref{app:corpus-selection} documents selection of these cases, and
Appendix~\ref{app:later-reference} gives a complete later-reference target.

\subsubsection{Feasibility and Scope}

I have applied the implemented single-concept operator to all 9,552 English
gold documents in the PMB training partition. Table~\ref{tab:coverage-funnel}
shows the resulting structural yield before the non-RSR hold-out.

\begin{table}[ht]
\centering
\small
\begin{tabular}{@{}lrr@{}}
\toprule
Stage & Count & \% of lexical concepts \\
\midrule
Lexical concepts (noun, verb, adjective) & 23,232 & 100.0 \\
Aligned to an uttered word               & 17,886 & 77.0 \\
Usable site with a WordNet reparandum    & 16,481 & 70.9 \\
Structurally legal transformation        & 16,425 & 70.7 \\
\bottomrule
\end{tabular}
\caption{Yield of the single-concept operator over the PMB training
partition. Of the usable sites, 99.7\% produce a legal transformation.}
\label{tab:coverage-funnel}
\end{table}

The 16,425 legal transformations form the candidate pool and include the
direct, anchor and movement constructions defined in
Section~\ref{sec:paired-configuration}. Structural generation and validation are complete;
linguistic filtering, the independent audit and final sampling remain. The
sampling rule limits repeated use of a source document and is fixed before
model training.

The main corpus covers immediate replacements of aligned noun, verb and
adjective concepts. It does not include role, tense, aspect or delayed
replacements. Its WordNet-derived reparanda and written PMB contexts also
differ from spontaneous speech. The experiments will therefore test learnability
within these boundaries, not coverage of naturally occurring RSR as a whole.


%==============New Model
\subsection{Model Experiment}
\label{sec:model}

The model experiment will test whether an LLM-based semantic parser can learn
the repair-aware representation from Section~\ref{sec:data}. A parse will count
as successful only if it recovers the withdrawn and accepted content, assigns
their status and links both to the surrounding clause. The evaluation will
also test later reference, accepted meaning and generalisation beyond the
training patterns.

\subsubsection{Experimental Design}

The experiment will use instruction-tuned Qwen3.5 models at three parameter
scales: 2B, 9B and 27B \citep{qwenQwen35Native2026}. Each model size will be
trained under all three conditions, allowing the effect of the training target
to be compared across model capacities.

All conditions will retain the fluent PMB pairs (Table~\ref{tab:model-experiment}).
RSR normalisation will add the generated utterances with fluent gold SBN targets that omit
the reparandum. Repair-aware training will use the same utterances with the
repair-aware targets from Section~\ref{sec:paired-configuration}.

\begin{table}[ht]
\centering
\small
\begin{tabular}{@{}p{0.22\textwidth}p{0.16\textwidth}p{0.28\textwidth}ccc@{}}
\toprule
Training condition & RSR inputs & Target for RSR inputs & 2B & 9B & 27B \\
\midrule
Fluent only
  & No  & ---                 & $\checkmark$ & $\checkmark$ & $\checkmark$ \\
RSR normalisation
  & Yes & Fluent gold SBN    & $\checkmark$ & $\checkmark$ & $\checkmark$ \\
Repair-aware
  & Yes & Repair-aware SBN    & $\checkmark$ & $\checkmark$ & $\checkmark$ \\
\bottomrule
\end{tabular}
\caption{Training conditions and model sizes. Every model size is trained
under all three conditions with three matched random seeds. The two RSR
conditions receive the same inputs and differ in the target assigned to each
RSR item.}
\label{tab:model-experiment}
\end{table}

Within each model size, the normalising and repair-aware conditions have their
inputs, example counts and training procedure held constant, so their
comparison isolates the effect of retaining withdrawn content and its ordered
replacement structure. The fluent-only condition provides a reference without
RSR supervision, but its comparison with repair-aware training combines new
examples and new targets. Results therefore pair gains in RSR representation
with changes in accepted-meaning and fluent-input accuracy.

All conditions will use the same parsing instruction and LoRA adaptation procedure
\citep{huLoRALowRankAdaptation2021}. The instruction does not identify an input
as fluent or RSR-containing. Each model--condition combination will use three
matched random seeds, giving 27 runs in total. Outputs will be evaluated
without automatic output correction. Appendix~\ref{app:model-config}
specifies the full procedure.

\subsubsection{Evaluation}

The primary measure is \emph{RSR-structure accuracy}. An item is correct
only if its output is a well-formed graph with the paired
\sbnterm{CORRECTION}--\sbnterm{CONJUNCTION} structure. It must also assign the
reparandum and repair to the withdrawn and accepted contexts and give both the
required links to shared content. Component scores show which requirement
failed.
Graph validity and failure-penalised SMATCH++ are complementary measures
of generation success and agreement with the complete target
\citep{opitzSMATCHStandardizedExtended2023}. Appendix~\ref{app:model-metrics}
gives the metric definitions and a worked scoring example.

Two controls assess how repair-aware training affects accepted-meaning
recovery and ordinary semantic parsing. For the
automatically generated RSR items, the accepted content of each prediction is
compared with the fluent gold PMB graph after RSR-specific structure is
removed. This measures recovery of the speaker's final intention independently
of success on the reparandum. The parsers are also evaluated on the unmodified
PMB test partition. Its graph score measures ordinary semantic parsing, while
the spurious-RSR rate records fluent inputs that are incorrectly assigned RSR
structure.

The later-reference challenge set tests the main functional motivation for
retaining withdrawn material. It checks links to the accepted referent, the
withdrawn referent and both referents where the target requires them. At each
model size, the normalising and repair-aware parsers are compared because only
the latter is trained to retain the withdrawn referent and its discourse
status.

Further tests examine whether the parser has learned an RSR distinction rather
than a narrow surface pattern. Paired versions of each test item measure the
effect of an editing expression while holding its lexical content constant, and
\textit{sorry} tests a marker excluded from training. The non-RSR set
measures false alarms on lists, appositions and coordination. Cross-turn
corrections test whether the within-clause configuration is applied outside
its intended scope.

The error diagnosis distinguishes failures of RSR structure, status and
semantic-role attachment from representations of editing expressions, spurious
RSR analyses and malformed SBN. Statistical comparisons use the original PMB
source document as the unit of observation, with its variants aggregated
before testing. For each model size, the three matched seeds are compared
separately and their mean difference and range are reported; a result is
treated as stable only when its direction is consistent across seeds.
Appendix~\ref{app:model-statistics} records the full error scheme and
statistical procedure.